\chapter{Обучение без «обучения»}
\label{ch:learning}

\section{Мы шлифуем веса, не занимаясь глубоким обучением}

Возникает верное ощущение: мы будто учим модель без «учить» и без «модели». Вроде
бы это не глубокое обучение и не графовая нейросеть, но тем не менее мы шлифуем
веса и параметры графа — по-своему. У этого ощущения есть точное содержание.

Система, которую мы строим, — параметрическая: у неё есть веса рёбер, уровни
уверенности, активации, и они обновляются опытом. Это \emph{обучение} в широком
смысле, даже без обратного распространения. Конкретно это два хорошо известных
паттерна, наложенных друг на друга:

\begin{itemize}
\item \textbf{Хеббовское/ассоциативное обучение} (семейство сетей Хопфилда;
принцип Хебба «что вместе активируется, вместе укрепляется»). Правило Хебба —
локальное усиление связи по со-активации:
\[
\Delta w_{ij}\;\propto\;a_i\,a_j ,
\]
а ассоциативная память Хопфилда хранит образы в минимумах энергии
$E=-\tfrac12\sum_{ij}w_{ij}a_i a_j$ и извлекает их спуском по $E$. Наши fusion и
консолидация, усиливающие со-активированные рёбра, — это буквально хеббовское
обучение на графе; граф памяти со spreading activation и уверенностью — современная
ассоциативная память, где извлечение есть спуск к аттрактору.
\item \textbf{Байесовское обновление уверенности} (глава~\ref{ch:propagation}):
локальные апдейты $\Delta\text{conf}$ со знаком, гейтом источника и дисконтом
независимости.
\end{itemize}

\section{Это message passing — но с локальными правилами}

Формально наша система — это \emph{передача сообщений} по графу (как в графовой
нейросети), но «параметры» (веса рёбер, уверенности) обновляются \emph{локальными}
правилами (хеббовскими, байесовскими, распадом), а не глобальным градиентным
спуском по функции потерь. Иначе говоря, это message passing \emph{без обратного
распространения}. Где мы сидим на карте методов — \emph{между} классической
ассоциативной памятью (Хопфилд) и графовыми нейросетями.

\begin{intuition}[Два способа настроить пианино]
Глубокое обучение настраивает все струны сразу по глобальному критерию («чтобы
аккорд звучал верно»), продвигаясь градиентом. Хеббовская система настраивает
каждую струну по локальному правилу («подтяни там, где чаще играют вместе»). Обе
«учат» инструмент, но первая знает про весь аккорд, а вторая — только про соседние
струны. Наша память — второго рода.
\end{intuition}

\section{Мост к «настоящей» графовой сети}

Может ли это стать графовой нейросетью в полном смысле? Да, и мост
поучителен: можно \emph{мета-обучать} сами правила обновления — учить оптимальное
правило $\Delta\text{conf}$, оптимальный распад, оптимальное смещение обхода
градиентным спуском вместо ручного дизайна. Это «обучение тому, как помнить»
(learning to remember), мета-обучение операторов памяти — реальное
исследовательское направление.

\begin{pitfall}[Но не на старте]
Мета-обучение правил меняет интерпретируемость на чёрный ящик, которому нужна
разметка и который заново вносит проблему непрозрачности субстрата. На старте —
ручные локальные правила: они интерпретируемы, отлаживаемы и не требуют обучающих
данных. «Может, и незачем» — верно для текущего этапа. Но мост реален, и полезно
знать, что стоишь ровно на нём: сегодняшняя система — предельный, «нулевой» случай
обучаемой графовой сети, где правила заданы руками, а не выучены.
\end{pitfall}

\section{Родство с DTI ещё на одном уровне}

Это же наблюдение ещё раз сшивает два проекта автора. Магистерская работа изучает
графовые нейросети с их обучаемыми операторами передачи сообщений; память Рин —
та же передача сообщений, но с локальными правилами. Если однажды правила памяти
начнут мета-обучаться, различие между «графовой моделью аффинности» и «графовой
моделью памяти» станет чисто предметным: один и тот же класс обучаемых операторов,
приложенный к разным графам. Это ещё один слой той единой спины, о которой шла речь
во введении. Существует и прямая работа на этом стыке — например, из библиографии
волта RIN: HeLa-Mem (Hebbian learning и ассоциативная память для LLM-агентов).
